\subsection{Flux vacua in type IIB compactifications on orbifolds: their finiteness and minimal string coupling --- arXiv:2404.18995}

\textbf{Authors:} Ignatios Antoniadis, Anthony Guillen, Osmin Lacombe

\paragraph{主な主張}
Type IIBトーラスorientifoldのフラックス真空を調べ、固定した\(N_{\rm flux}\)で非同値な真空が有限であることと、十分大きな電荷で\(g_{s,\min}\sim c/N_{\rm flux}^{\alpha}\), \(\alpha=1,2\)となることに強い数値的根拠を与えた。 \cite{Antoniadis:2024hqd}

\paragraph{新規性}
真空の数え上げと弦結合の下限を結び付け、磁化D7ブレーンが誘起する負のD3電荷によりtadpole上限が緩和され、より弱い結合へ到達できることを示した。

\paragraph{位置付け}
弱結合による有効理論の制御と、大きなフラックスに対する電荷制約の両立を定量的に検討する研究である。

\paragraph{コメント（読書メモ）}
檜垣さんからのコメント（原メモを整文）：Type IIBのトーラスorbifold上のorientifoldで、フラックスによるモジュライ固定が議論されている。弦結合の最小値は\(1/N\)または\(1/N^2\)に比例する傾向があり、\(N\)（フラックスが運ぶRR電荷）が大きいほど弱結合になる。反Dブレーン電荷と超対称性破れがあると制約が緩み、さらに弱結合となる傾向がある。一方、大きな\(N\)は摂動論の条件を満たしやすくするが、tadpole条件（元メモではanomalyに関する制限との類似を指摘）を満たしにくくするため、「こちらを立てるとあちらが立たない」状況と感じる。
